\documentclass[aspectratio=169,11pt]{beamer}
\usepackage[T1]{fontenc}
\usepackage[utf8]{inputenc}
\usepackage[french]{babel}
\usepackage{lmodern}
\usepackage{booktabs}
\usepackage{tabularx}
\usepackage{tikz}
\usepackage{fontawesome5}
\usetikzlibrary{positioning,arrows.meta,calc,fit}

\usetheme[progressbar=frametitle,numbering=fraction]{metropolis}

\definecolor{nuit}{HTML}{14213D}
\definecolor{papier}{HTML}{F7F6F2}
\definecolor{bleu}{HTML}{1F5FAD}
\definecolor{bleuclair}{HTML}{DCE7F6}
\definecolor{orange}{HTML}{C2581F}
\definecolor{orangeclair}{HTML}{F6E6DA}
\definecolor{vert}{HTML}{2E7D32}
\definecolor{gris}{HTML}{4A5568}
\definecolor{trait}{HTML}{DDD8CC}

\setbeamercolor{normal text}{fg=nuit,bg=papier}
\setbeamercolor{frametitle}{fg=papier,bg=nuit}
\setbeamercolor{progress bar}{fg=bleu,bg=bleuclair}
\setbeamercolor{alerted text}{fg=orange}
\setbeamercolor{title separator}{fg=bleu}
\setbeamercolor{background canvas}{bg=papier}

\tikzset{
  carte/.style={draw=trait,fill=white,rounded corners=4pt,align=center,inner sep=6pt,font=\small,execute at begin node={\hyphenpenalty=10000}},
  fleche/.style={-{Stealth[length=2.5mm]},thick,bleu},
}


% Probabilités de l'exemple (grille 5 x 5) : colonne/ligne/probabilité, ligne 1 en haut
\newcommand{\cellules}{2/2/0.05, 3/2/0.10, 4/2/0.05, 1/3/0.05, 2/3/0.15, 3/3/0.30, 4/3/0.10, 2/4/0.05, 3/4/0.10, 4/4/0.05}
% Grille vide 5 x 5, coin haut gauche en (0,0), cases de 0.75
\newcommand{\grille}{%
  \foreach \i in {0,...,4} \foreach \j in {0,...,4} \draw[trait,fill=white] (\i*0.75,-\j*0.75) rectangle ++(0.75,-0.75);}
% Cases remplies selon leur probabilité, avec la valeur écrite
\newcommand{\carte}{%
  \foreach \c/\l/\p in \cellules {
    \pgfmathsetmacro{\op}{\p/0.30}
    \fill[orange,fill opacity=\op] (\c*0.75,-\l*0.75+0.75) rectangle ++(0.75,-0.75);
    \node[font=\tiny] at (\c*0.75+0.375,-\l*0.75+0.375) {\p};}}

\title{La tête séparable}
\subtitle{Comment le modèle transforme un nœud en position, et comment le faire 250 fois moins cher}
\author{Melissa Colin}
\date{}

\begin{document}

\maketitle

%---------------------------------------------------------------------
\begin{frame}{Le rôle de la tête}
\centering
\begin{tikzpicture}[node distance=0.6cm]
  \node[carte] (v) {\textbf{Un nœud}\\\scriptsize vecteur de 64 nombres\\\scriptsize calculé par le GNN};
  \node[carte,right=of v,fill=bleuclair,draw=bleu] (t) {\textbf{La tête}\\\scriptsize dernière couche};
  \node[carte,right=of t] (p) {\textbf{Sa position}\\\scriptsize $(x, y)$ en pixels};
  \draw[fleche] (v) -- (t);
  \draw[fleche] (t) -- (p);
  \begin{scope}[shift={($(p.east)+(0.8,1.1)$)}]
    \draw[trait,fill=white] (0,0) rectangle (2.2,-2.2);
    \fill[bleu] (1.5,-0.9) circle (2.5pt);
    \draw[gris,dashed] (1.5,0) -- (1.5,-0.9) -- (0,-0.9);
    \node[font=\tiny,gris,above] at (1.5,0) {$x$};
    \node[font=\tiny,gris,left] at (0,-0.9) {$y$};
    \node[font=\tiny,gris,below] at (1.1,-2.2) {image $500 \times 500$};
  \end{scope}
\end{tikzpicture}

\vspace{0.8cm}
\small La question de ces slides : \textbf{comment la tête passe du vecteur à la position}, et pourquoi on peut le faire beaucoup moins cher.

\vspace{0.2cm}
\scriptsize\textcolor{gris}{Pour que ce soit lisible, l'exemple utilise une image de $5 \times 5$ pixels au lieu de $500 \times 500$. Le principe est le même.}
\end{frame}

%---------------------------------------------------------------------
\begin{frame}{Tête complète : une note pour chaque pixel}
\begin{columns}[c]
\begin{column}{0.42\textwidth}
\centering
\begin{tikzpicture}
  \grille \carte
  \foreach \i in {1,...,5} {
    \node[font=\scriptsize,gris] at (\i*0.75-0.375,0.3) {\i};
    \node[font=\scriptsize,gris] at (-0.3,-\i*0.75+0.375) {\i};}
  \node[font=\scriptsize,gris] at (1.875,0.7) {colonne $x$};
  \node[font=\scriptsize,gris,rotate=90] at (-0.75,-1.875) {ligne $y$};
\end{tikzpicture}
\end{column}
\begin{column}{0.56\textwidth}
\small
\begin{enumerate}
\setlength{\itemsep}{6pt}
\item La tête donne \textbf{une note à chaque pixel} : 25 notes ici, \textbf{250\,000} pour $500 \times 500$.
\item Un softmax les change en \textbf{probabilités} (total = 1) : une carte « où le nœud veut être ».
\item La position est la \textbf{moyenne pondérée} :\\[3pt]
$x = \sum_{\text{pixels}} p \times \text{colonne} = 3{,}85$\\[2pt]
$y = \sum_{\text{pixels}} p \times \text{ligne} = 3{,}00$
\end{enumerate}

\vspace{0.2cm}
\scriptsize\textcolor{gris}{250\,000 notes à partir de 64 nombres : $64 \times 250\,000 \approx 16$ millions de poids, 95\,\% du modèle.}
\end{column}
\end{columns}
\end{frame}

%---------------------------------------------------------------------
\begin{frame}{Ce que la moyenne utilise vraiment}
\begin{columns}[c]
\begin{column}{0.48\textwidth}
\centering
\begin{tikzpicture}
  \grille \carte
  % totaux par colonne, sous la grille
  \foreach \c/\t in {1/0, 2/0.05, 3/0.25, 4/0.50, 5/0.20} {
    \pgfmathsetmacro{\op}{\t/0.5}
    \fill[bleu,fill opacity=\op] (\c*0.75-0.75,-4.05) rectangle ++(0.75,-0.55);
    \draw[bleu] (\c*0.75-0.75,-4.05) rectangle ++(0.75,-0.55);
    \node[font=\tiny] at (\c*0.75-0.375,-4.325) {\t};}
  % totaux par ligne, à droite
  \foreach \l/\t in {1/0, 2/0.20, 3/0.60, 4/0.20, 5/0} {
    \pgfmathsetmacro{\op}{\t/0.6}
    \fill[vert,fill opacity=\op] (4.05,-\l*0.75+0.75) rectangle ++(0.55,-0.75);
    \draw[vert] (4.05,-\l*0.75+0.75) rectangle ++(0.55,-0.75);
    \node[font=\tiny] at (4.325,-\l*0.75+0.375) {\t};}
  \draw[-{Stealth},bleu,thick] (1.875,-3.8) -- (1.875,-4.0);
  \node[font=\scriptsize,bleu] at (1.875,-4.95) {total de chaque colonne};
  \node[font=\scriptsize,vert,rotate=-90] at (5.0,-1.875) {total de chaque ligne};
\end{tikzpicture}
\end{column}
\begin{column}{0.50\textwidth}
\small
On regroupe les termes de la moyenne :\\[4pt]
$x = 1 \times \textcolor{bleu}{0} + 2 \times \textcolor{bleu}{0{,}05} + 3 \times \textcolor{bleu}{0{,}25} + 4 \times \textcolor{bleu}{0{,}50} + 5 \times \textcolor{bleu}{0{,}20} = 3{,}85$\\[3pt]
$y = 2 \times \textcolor{vert}{0{,}20} + 3 \times \textcolor{vert}{0{,}60} + 4 \times \textcolor{vert}{0{,}20} = 3{,}00$

\vspace{0.4cm}
\begin{tikzpicture}
  \node[carte,draw=orange,fill=orangeclair,text width=0.9\linewidth,align=left] {\small $x$ ne dépend que des \textcolor{bleu}{\textbf{totaux par colonne}}, $y$ que des \textcolor{vert}{\textbf{totaux par ligne}}. Le détail de la carte, lui, ne change pas la position.};
\end{tikzpicture}
\end{column}
\end{columns}
\end{frame}

%---------------------------------------------------------------------
\begin{frame}{Tête séparable : prédire directement les totaux}
\begin{columns}[c]
\begin{column}{0.46\textwidth}
\centering
\begin{tikzpicture}
  \node[font=\small\bfseries] at (1.875,0.5) {5 notes pour les colonnes};
  \foreach \c/\t in {1/0, 2/0.05, 3/0.25, 4/0.50, 5/0.20} {
    \pgfmathsetmacro{\op}{\t/0.5}
    \fill[bleu,fill opacity=\op] (\c*0.75-0.75,0) rectangle ++(0.75,-0.55);
    \draw[bleu] (\c*0.75-0.75,0) rectangle ++(0.75,-0.55);
    \node[font=\tiny] at (\c*0.75-0.375,-0.275) {\t};}
  \node[font=\small,bleu] at (1.875,-1.0) {$\Rightarrow x = 3{,}85$};
  \node[font=\small\bfseries] at (1.875,-1.9) {5 notes pour les lignes};
  \foreach \l/\t in {1/0, 2/0.20, 3/0.60, 4/0.20, 5/0} {
    \pgfmathsetmacro{\op}{\t/0.6}
    \fill[vert,fill opacity=\op] (\l*0.75-0.75,-2.4) rectangle ++(0.75,-0.55);
    \draw[vert] (\l*0.75-0.75,-2.4) rectangle ++(0.75,-0.55);
    \node[font=\tiny] at (\l*0.75-0.375,-2.675) {\t};}
  \node[font=\small,vert] at (1.875,-3.4) {$\Rightarrow y = 3{,}00$};
\end{tikzpicture}
\end{column}
\begin{column}{0.52\textwidth}
\small
\begin{itemize}
\setlength{\itemsep}{6pt}
\item Au lieu de la carte entière, la tête prédit \textbf{une note par colonne} et \textbf{une par ligne} : $5 + 5 = 10$ notes ici, \textbf{1\,000} pour $500 \times 500$.
\item Un softmax pour les colonnes, un pour les lignes, puis la même moyenne pondérée : \textbf{même position} $(3{,}85\,;\,3{,}00)$.
\item \textbf{Rien n'est perdu} : toute position de la tête complète s'obtient ainsi, et inversement.
\end{itemize}
\end{column}
\end{columns}
\end{frame}

%---------------------------------------------------------------------
\begin{frame}{Ce qu'on y gagne, pour une image de $500 \times 500$}
\centering
\small
\begin{tabular}{@{}lcc@{}}
\toprule
 & \textbf{Tête complète} & \textbf{Tête séparable} \\
\midrule
Stress normalisé $\downarrow$ & 0,065 & \textbf{0,064} \\
Croisements d'arêtes $\downarrow$ & 30,5 & \textbf{30,4} \\
Nœuds superposés $\downarrow$ & 0,06\,\% & \textbf{0,04\,\%} \\
\midrule
Notes par nœud & 250\,000 & \textbf{1\,000} \\
Paramètres du modèle entier & 17\,M & \textbf{841\,k} \quad\scriptsize\textcolor{gris}{(CoRe-GD : 930\,k)} \\
Mémoire GPU pendant le dessin & 289 Mo & \textbf{3 Mo} \\
Temps par graphe & 18 ms & \textbf{12 ms} \\
Une epoch d'entraînement & $\sim$3 min & \textbf{$\sim$1 min 15} \\
Epochs (même temps de calcul) & 15 & 30 au plus \\
\bottomrule
\end{tabular}

\vspace{0.25cm}
\scriptsize\textcolor{gris}{Les 1\,000 graphes de test de Rome, même loss (stress seul). La tête séparable apprend plus lentement par epoch, mais chaque epoch coûte 2,6 fois moins : à temps de calcul égal, elle fait aussi bien ou mieux partout.}
\end{frame}

\end{document}
